\chapter{Generating Functions and the Symbolic Method}

\begin{center}\itshape
Whatever pattern of three coin tosses you choose, I can choose one that is more likely to appear first.\\
\upshape(the moral of Walter Penney's game, 1969, in our words)
\end{center}

Here is a bet to offer a friend. Each of you picks a pattern of three coin tosses, say HHT for your friend; then a coin is
tossed until one of the two patterns appears, and its owner wins. Your friend chooses first, and you always win in the
long run. Against HHT you choose THH, and win three times out of four. Against THH you choose TTH and win two times out of
three; against TTH, HTT, which wins three times out of four; and against HTT, HHT, which wins two times out of three. The
four patterns beat each other in a circle, like rock, paper and scissors (Figure~\ref{fig:penney}). Walter Penney
published the game in 1969 in the \emph{Journal of Recreational Mathematics}; John Conway found a strange rule that
computes the odds for any two patterns, and Martin Gardner printed it in his \emph{Scientific American} column in 1974, without a
proof. The proof came in 1981, from Leo Guibas and Andrew Odlyzko, and it is a short computation with generating functions.

\begin{figure}[ht]
\centering
\fig{ch04-penney}
\caption{Penney's game. An arrow from $A$ to $B$ means that $B$ appears before $A$ with the odds shown: each pattern of the
cycle is beaten by the next one. Whatever the first player chooses, the second player has a better pattern
(Exercise~\ref{ex:penney}).}
\label{fig:penney}
\end{figure}

Tome~4, Chapter~2 used generating functions as a computing device: a sequence becomes a power series, and awkward
operations on sequences become easy operations on functions. This chapter adds the idea that turned them into a
method, developed by Philippe Flajolet and his school from the 1980s: the \emph{symbolic method}. Describe a family of
objects by the way it is built (``a tree is a root and a sequence of trees'', ``a word avoiding a pattern is\dots''), and a
dictionary translates the description into an equation for its generating function, with no computation at all. Then
algebra extracts exact counts, and analysis (Tome~4, Chapter~11) their asymptotic size. It is the tool behind the
average-case analysis of algorithms, pioneered by Donald Knuth in \emph{The Art of Computer Programming}: how deep is a
random tree, how long until a search succeeds, how many cycles does a random permutation have. We will also meet the
Lagrange inversion formula, which counts trees of every kind at once, and prove it with the cycle lemma of Chapter~3.

\begin{think}
Toss a fair coin until HH appears, and separately until HT appears. Both patterns have probability $\frac14$ at each position.
Do they take the same time on average? If not, which one takes longer, and why?
\end{think}

\section{The symbolic method}

\begin{reminder}[Tome~4, Chapter~2]
Formal power series over $\Q$ form a ring $\Q[[z]]$. A series is invertible if and only if its constant term is non-zero
(Proposition~2.1). Sums $\sum_kA_k$ and products $\prod_k(1+A_k)$ make sense when the orders of the $A_k$ tend to infinity, and so does
the composition $F(G)$ when $G(0)=0$. The derivative obeys the product and chain rules, and $\exp$ and $\log$ satisfy
$\exp(F+G)=\exp F\exp G$ and $\exp(\log(1+F))=1+F$ when $F(0)=G(0)=0$ (Proposition~2.2).
\end{reminder}

A \emph{combinatorial class} is a set $\mathcal A$ of objects with a \emph{size} $|\alpha|\in\N$ for each object, such that there are finitely many
objects of each size. Its \emph{counting sequence} is $a_n=|\{\alpha:|\alpha|=n\}|$ and its \emph{ordinary generating function} (OGF) is
\[
A(z)=\sum_{\alpha\in\mathcal A}z^{|\alpha|}=\sum_{n\ge0}a_nz^n .
\]
The basic classes are $\mathcal E$, one object of size $0$, with OGF $1$, and $\mathcal Z$, one object of size $1$ (an \emph{atom}), with OGF $z$. New
classes are built from old ones by constructions:
\begin{itemize}
\item the \emph{disjoint union} $\mathcal A+\mathcal B$, keeping the sizes;
\item the \emph{product} $\mathcal A\times\mathcal B$, the pairs $(\alpha,\beta)$ with size $|\alpha|+|\beta|$;
\item the \emph{sequences} $\operatorname{SEQ}(\mathcal A)$, the finite tuples $(\alpha_1,\dots,\alpha_k)$, $k\ge0$, with size $|\alpha_1|+\dots+|\alpha_k|$ (we require $a_0=0$, otherwise
there would be infinitely many sequences of size $0$);
\item the \emph{multisets} $\operatorname{MSET}(\mathcal A)$ and the \emph{sets} $\operatorname{PSET}(\mathcal A)$: sequences up to order, with repetitions allowed or not ($a_0=0$);
\item the \emph{cycles} $\operatorname{CYC}(\mathcal A)$: sequences of length $k\ge1$ up to cyclic rotation ($a_0=0$).
\end{itemize}

\begin{theorem}[the symbolic dictionary]\label{thm:dictionary}
With $A$, $B$ the OGFs of $\mathcal A$, $\mathcal B$:
\[
\begin{aligned}
\mathcal A+\mathcal B&\mapsto A(z)+B(z),\\
\mathcal A\times\mathcal B&\mapsto A(z)B(z),\\
\operatorname{SEQ}(\mathcal A)&\mapsto\frac1{1-A(z)},\\
\operatorname{MSET}(\mathcal A)&\mapsto\prod_{n\ge1}(1-z^n)^{-a_n}\eqbr=\exp\Big(\sum_{k\ge1}\frac{A(z^k)}k\Big),\\
\operatorname{PSET}(\mathcal A)&\mapsto\prod_{n\ge1}(1+z^n)^{a_n}\eqbr=\exp\Big(\sum_{k\ge1}\frac{(-1)^{k-1}A(z^k)}k\Big),\\
\operatorname{CYC}(\mathcal A)&\mapsto\sum_{d\ge1}\frac{\varphi(d)}d\log\frac1{1-A(z^d)} .
\end{aligned}
\]
\end{theorem}
\begin{proof}
Union: the sum over $\mathcal A+\mathcal B$ of $z^{|\gamma|}$ splits into two sums. Product: $\sum_{(\alpha,\beta)}z^{|\alpha|+|\beta|}=\sum_\alpha z^{|\alpha|}\sum_\beta z^{|\beta|}$, a product of series
(each coefficient is a finite sum). Sequences: by the product rule, the sequences of length $k$ have OGF $A^k$, of order at least $k$
since $a_0=0$; summing over $k$, $\sum_kA^k=\frac1{1-A}$, the inverse of $1-A$ (Tome~4, Proposition~2.1).

Multisets and sets: a multiset is the choice of a multiplicity $m_\alpha\ge0$ for each object, finitely many non-zero, and contributes
$\prod_\alpha z^{m_\alpha|\alpha|}$; the choices are independent, so the OGF is $\prod_\alpha(1+z^{|\alpha|}+z^{2|\alpha|}+\cdots)=\prod_n(1-z^n)^{-a_n}$, with $m_\alpha\in\{0,1\}$
for sets. The coefficient of $z^N$ only involves the factors with $n\le N$, so the infinite products make sense. For the exponential
forms, $\frac1{1-z^n}=\exp\big(\sum_kz^{nk}/k\big)$ and $1+z^n=\exp\big(\sum_k(-1)^{k-1}z^{nk}/k\big)$ by Proposition~2.2, and multiplying these exponentials
adds their arguments: $\sum_na_n\sum_kz^{nk}/k=\sum_kA(z^k)/k$, and similarly with signs.

Cycles: the cycles of length $k$ are the orbits of the cyclic group of rotations acting on $\mathcal A^k$. By P\'olya's theorem
(Tome~2, Theorem~3.7), with the objects as colours weighted by $z^{|\alpha|}$, their OGF is the cycle index of the rotations of $k$
slots evaluated at $t_j=A(z^j)$. A rotation by $r$ steps has $\gcd(r,k)$ cycles of length $d=k/\gcd(r,k)$, and $\varphi(d)$ rotations have this
$d$ (Tome~2, Example~3.8), so the OGF is $\frac1k\sum_{d\mid k}\varphi(d)A(z^d)^{k/d}$. (There are infinitely many colours, but the coefficient of $z^N$
only involves objects of size $\le N$, as in Proposition~\ref{prop:multisets}.) Summing over $k=dm$,
\[
\begin{aligned}{}&\sum_{k\ge1}\frac1k\sum_{d\mid k}\varphi(d)A(z^d)^{k/d}=\sum_{d\ge1}\frac{\varphi(d)}d\sum_{m\ge1}\frac{A(z^d)^m}m\eqbr=\sum_{d\ge1}\frac{\varphi(d)}d\log\frac1{1-A(z^d)}.\end{aligned}
\]
\end{proof}

The point is that one never computes: one \emph{describes}, and the theorem computes. Some first translations:
\begin{itemize}
\item binary words, $\operatorname{SEQ}(\mathcal Z+\mathcal Z)$: $\frac1{1-2z}$, and $2^n$ words of length $n$;
\item compositions of $n$ (ordered sums of positive integers), $\operatorname{SEQ}(\operatorname{SEQ}_{\ge1}(\mathcal Z))$: $\frac1{1-z/(1-z)}=\frac{1-z}{1-2z}$, so $2^{n-1}$ of them for $n\ge1$;
\item partitions of $n$, $\operatorname{MSET}(\operatorname{SEQ}_{\ge1}(\mathcal Z))$: Euler's $\prod_n\frac1{1-z^n}$ (Tome~4, Proposition~2.11);
\item binary necklaces, $\operatorname{CYC}(\mathcal Z+\mathcal Z)$: the formula of Tome~2, Example~3.8 (Exercise~\ref{ex:necklacecyc});
\item plane trees, $\mathcal G=\mathcal Z\times\operatorname{SEQ}(\mathcal G)$: $G=\frac z{1-G}$, the Catalan equation of Tome~3, Proposition~3.10;
\item rooted unlabelled trees, $\mathcal T=\mathcal Z\times\operatorname{MSET}(\mathcal T)$: Cayley's equation of Chapter~3.
\end{itemize}
A recursive specification such as $\mathcal G=\mathcal Z\times\operatorname{SEQ}(\mathcal G)$ determines the counting sequence, since the objects of size $n$ are built from
smaller ones; it is a grammar, and its generating function is the solution of the corresponding equation.

\subsection*{Automata}

A whole family of specifications needs only union and product, and gives rational generating functions. A
\emph{deterministic finite automaton} over an alphabet $\Sigma$ has a finite set of states $Q$, an initial state $s_0$, a set $F\subseteq Q$ of accepting
states, and a transition function $\delta:Q\times\Sigma\to Q$. It reads a word letter by letter, moving from state to state, and \emph{accepts} the word if it
stops in $F$. The \emph{transition matrix} $M$ has entry $M_{st}$ equal to the number of letters $a$ with $\delta(s,a)=t$.

\begin{theorem}\label{thm:automaton}
The OGF of the words accepted by an automaton is $L(z)=e_{s_0}^{\mathsf T}(I-zM)^{-1}\mathbf 1_F$, where $\mathbf 1_F$ is the indicator vector of $F$. In particular it is a
rational function of $z$, and the number of accepted words of length $n$ satisfies a linear recurrence with constant coefficients.
\end{theorem}
\begin{proof}
Let $L_s(z)$ be the OGF of the words that lead from state $s$ to an accepting state. A word is either empty, accepted from $s$ if $s\in F$, or
a letter $a$ followed by a word accepted from $\delta(s,a)$. So $L_s=[s\in F]+z\sum_aL_{\delta(s,a)}=[s\in F]+z\sum_tM_{st}L_t$; in vector form $L=\mathbf 1_F+zML$, that
is, $(I-zM)L=\mathbf 1_F$. The determinant of $I-zM$ is a polynomial with constant term $1$, invertible in $\Q[[z]]$, so by Cramer's rule
$L=(I-zM)^{-1}\mathbf 1_F$, with rational entries. A rational OGF $\frac PQ$ with $Q(0)=1$ has coefficients satisfying the recurrence given by $Q$
(Tome~4, Theorem~2.5).
\end{proof}

\begin{figure}[ht]
\centering
\fig{ch04-automaton}
\caption{An automaton accepting the binary words without two consecutive $1$s. State $a$: the last letter is not a $1$;
state $b$: it is a $1$. Both states accept; reading $1$ in state $b$ leads to a rejecting sink, not drawn. Its transition
matrix restricted to $\{a,b\}$ is $\left(\begin{smallmatrix}1&1\\1&0\end{smallmatrix}\right)$, whose characteristic polynomial gives the Fibonacci recurrence.}
\label{fig:automaton}
\end{figure}

For the words with no two consecutive $1$s (Figure~\ref{fig:automaton}), $L_a=1+zL_a+zL_b$ and $L_b=1+zL_a$, so $L_a=\frac{1+z}{1-z-z^2}$, and there
are $F_{n+2}$ such words of length $n$. Regular expressions, lexers and network protocols are automata, and this theorem is how one
counts what they accept. Chomsky and Sch\"utzenberger (1963) extended it one level up: an unambiguous context-free grammar has an
\emph{algebraic} generating function, like the Dyck language with its Catalan equation.

\section{Patterns in words}

Now Penney's game. We work with an alphabet of $q$ letters and with \emph{patterns}, words $A$, $B$ of the same length $k$. An \emph{occurrence} of $A$ in a
word is a factor (a block of consecutive letters) equal to $A$; it \emph{ends} at the position of its last letter. Occurrences can overlap:
in HHHH, the pattern HH occurs three times. Overlaps are measured by the \emph{correlation}: for words $X$, $Y$ of length $k$, let
$O(X,Y)$ be the set of $j\in\{1,\dots,k\}$ such that the last $j$ letters of $X$ are the first $j$ letters of $Y$, and define
\[
\begin{gathered}c_{XY}(z)=\sum_{j\in O(X,Y)}z^{k-j},\qq XY=\sum_{j\in O(X,Y)}q^j=q^kc_{XY}(1/q).\end{gathered}
\]
For instance, with $q=2$: $O(\mathrm{HH},\mathrm{HH})=\{1,2\}$, so $\mathrm{HH\,HH}=2+4=6$, while $O(\mathrm{HT},\mathrm{HT})=\{2\}$ and $\mathrm{HT\,HT}=4$. Conway wrote the binary number
whose $j$-th digit from the right says whether $j\in O(X,Y)$, and called it the \emph{leading number} of $X$ over $Y$; $XY$ is its value. The set $O(X,X)$
is exactly what the Knuth--Morris--Pratt string-matching algorithm precomputes (Tome~C2).

\begin{theorem}[Guibas and Odlyzko, 1981]\label{thm:guibas}
Let $A\ne B$ be words of length $k$ over $q$ letters. Let $\mathcal S$ be the class of words containing neither $A$ nor $B$, and $\mathcal T_A$ (resp.\ $\mathcal T_B$) the
class of words that end with an occurrence of $A$ (resp.\ $B$) and contain no other occurrence of $A$ or $B$. Then
\[
\begin{aligned}1+qz\,S(z)&=S(z)+T_A(z)+T_B(z),\\ z^kS(z)&=c_{AA}(z)T_A(z)+c_{BA}(z)T_B(z),\\ z^kS(z)&=c_{AB}(z)T_A(z)+c_{BB}(z)T_B(z).\end{aligned}
\]
With a single pattern $A$ (and $\mathcal T_B$ empty), $1+qzS=S+T_A$ and $z^kS=c_{AA}T_A$.
\end{theorem}
\begin{proof}
\emph{First identity.} Adding a letter to a word of $\mathcal S$ (in $q$ ways) gives a non-empty word that either still contains neither pattern, or
contains one for the first time, at its end: a word of $\mathcal S$, $\mathcal T_A$ or $\mathcal T_B$. Conversely every non-empty word of these classes arises exactly
once, by deleting its last letter, which leaves a word containing neither pattern. With the empty word, $1+qzS=S+T_A+T_B$.

\emph{Second identity.} Take $w\in\mathcal S$ and form $wA$. It contains an occurrence of $A$ or $B$, and all of them end in the appended $A$. Let the first
one end $j$ letters into the appended $A$, $1\le j\le k$. If it is an occurrence of $A$, its last $j$ letters are the first $j$ letters of $A$, so $j\in O(A,A)$,
and $wA=ur$ where $u$, the prefix of $wA$ ending there, lies in $\mathcal T_A$, and $r$ is the last $k-j$ letters of $A$. If it is an occurrence of $B$, likewise
$j\in O(B,A)$ (note $j<k$, since $B\ne A$) and $u\in\mathcal T_B$. Conversely, given $u\in\mathcal T_A$ and $j\in O(A,A)$, delete the last $j$ letters of $u$ to get $w$; then $w$
contains neither pattern (it is a proper prefix of $u$), $wA=ur$, and the first occurrence of a pattern in $wA$ is the one ending at the end of $u$,
since an earlier one would lie in $u$. The same holds for $u\in\mathcal T_B$ and $j\in O(B,A)$. This bijection $w\leftrightarrow(u,j)$ satisfies $|w|+k=|u|+(k-j)$, which
gives $z^kS=c_{AA}T_A+c_{BA}T_B$. The third identity is the same with $B$ appended. With a single pattern, drop $B$ everywhere.
\end{proof}

Now flip a fair coin forever, or more generally draw independent uniform letters. The probability that the first $n$ letters form a given
word is $q^{-n}$, so substituting $z=1/q$ turns these counting series into probabilities. Let $\tau$ be the first time one of the
patterns is complete.

\begin{corollary}\label{cor:penney}
\begin{enumerate}
\item The words of length $n$ avoiding a pattern $A$ of length $k$ have OGF $\displaystyle\frac{c_{AA}(z)}{z^k+(1-qz)c_{AA}(z)}$.
\item The expected waiting time for $A$ is $\mathbb E[\tau]=AA$.
\item (Conway's formula) $A$ appears before $B$ with odds $\mathbb P(A\text{ first}):\mathbb P(B\text{ first})=(BB-BA):(AA-AB)$.
\end{enumerate}
\end{corollary}
\begin{proof}
(1) Eliminate $T_A=z^kS/c_{AA}$ from $1+qzS=S+T_A$ (note $c_{AA}(0)=1$, as $k\in O(A,A)$, so $c_{AA}$ is invertible).

For (2) and (3), all the series have non-negative coefficients, and $q^{-n}[z^n]S=\mathbb P(\tau>n)$, while $q^{-n}[z^n]T_A=\mathbb P(\tau=n$ and $A$ completes first$)$. Cutting the
letters into blocks of $k$, each block is $A$ with probability $q^{-k}$ independently, so $\mathbb P(\tau>n)\le(1-q^{-k})^{\lfloor n/k\rfloor}$, and $S(1/q)=\sum_n\mathbb P(\tau>n)=\mathbb E[\tau]$
converges (Tome~3, Chapter~2); $T_A(1/q)=\mathbb P(A\text{ first})$. The identities of Theorem~\ref{thm:guibas} relate convergent series with non-negative
coefficients, so they hold at $z=1/q$, where the first one reads $1+S=S+T_A+T_B$: one of the patterns appears, with probability $1$.

(2) With a single pattern, $T_A(1/q)=1$, and $q^{-k}S(1/q)=c_{AA}(1/q)$, so $\mathbb E[\tau]=q^kc_{AA}(1/q)=AA$.

(3) Multiply the second and third identities at $z=1/q$ by $q^k$: $\mathbb E[\tau]=AA\,p_A+BA\,p_B=AB\,p_A+BB\,p_B$, where $p_A=\mathbb P(A\text{ first})$. Subtracting,
$p_A(AA-AB)=p_B(BB-BA)$.
\end{proof}

So the answer to the opening question: HH takes $6$ tosses on average and HT only $4$. An occurrence of HT cannot overlap the next one, while
occurrences of HH come in clumps (HHH contains two); the long-run frequency is the same, $\frac14$, so the clumps must be further
apart. Tome~3, Chapter~10 proved (2) with a martingale and a casino of gamblers betting on ABRACADABRA, which a monkey types after
$26^{11}+26^4+26$ letters on average; here it is one line of algebra.

For Penney's game, take $A=\mathrm{THH}$ and $B=\mathrm{HHT}$. Then $AA=8$ (only $j=3$), $AB=2+4=6$ (THH ends with H and with HH, and HHT starts with
both), $BB=8$, and $BA=2$ (HHT ends with T, and THH starts with T). The odds of THH are $(BB-BA):(AA-AB)=(8-2):(8-6)=3:1$. The second player's rule is simple: against $b_1b_2b_3$, play $\bar b_2b_1b_2$, where $\bar b$ is the other side of the coin
(Exercise~\ref{ex:penney}).

\begin{beauty}
Three linear equations, read off a picture of overlapping words, contain the number of words avoiding a pattern, the expected
waiting time, Conway's mysterious rule, and the non-transitivity of the game. The paradox dissolves: a pattern that overlaps
itself, like HHH, wastes its chances in clumps, while a pattern that overlaps the beginning of its opponent, like THH against HHH,
lies in wait just before it.
\end{beauty}

\section{Labelled structures}

Many structures live on a set of labels: permutations, set partitions, graphs, maps. For them the \emph{exponential} generating
function $\hat A(z)=\sum_na_n\frac{z^n}{n!}$ is the right tool, as Tome~4, Chapter~2 explained. A \emph{labelled class} is a class whose objects of size $n$
carry the labels $1,\dots,n$, one on each atom. The \emph{labelled product} $\mathcal A\star\mathcal B$ consists of the pairs $(\alpha,\beta)$ relabelled in all possible ways that keep the
relative order of the labels inside $\alpha$ and inside $\beta$: choose which $|\alpha|$ of the labels $1,\dots,|\alpha|+|\beta|$ go to $\alpha$. It is the reason for the binomial
coefficient in $\hat A\hat B=\sum_n\big(\sum_k\binom nka_kb_{n-k}\big)\frac{z^n}{n!}$. Iterating, $\operatorname{SEQ}_k(\mathcal A)=\mathcal A^{\star k}$; and since the components of an element of $\mathcal A^{\star k}$ carry
disjoint sets of labels, they are distinct, so the $k!$ reorderings of a $k$-tuple are distinct tuples, and so are its $k$ rotations.

\begin{theorem}[labelled dictionary]\label{thm:labelled}
For labelled classes with $a_0=0$:
\[
\begin{gathered}
\mathcal A\star\mathcal B\mapsto\hat A\hat B,\qquad \operatorname{SEQ}(\mathcal A)\mapsto\frac1{1-\hat A},\qq \operatorname{SET}(\mathcal A)\mapsto\exp\hat A,\qquad \operatorname{CYC}(\mathcal A)\mapsto\log\frac1{1-\hat A}.
\end{gathered}
\]
\end{theorem}
\begin{proof}
The product is the binomial convolution above, and $\operatorname{SEQ}_k$ gives $\hat A^k$. A set of $k$ components corresponds to exactly $k!$ sequences, its
orderings, so $\operatorname{SET}_k\mapsto\frac{\hat A^k}{k!}$; a cycle of $k$ components to exactly $k$ sequences, its rotations, so $\operatorname{CYC}_k\mapsto\frac{\hat A^k}k$. Sum over $k$.
\end{proof}

The rule for sets is the exponential formula of Tome~4, Theorem~2.9. Here is what the dictionary says at once:
\begin{itemize}
\item permutations are sets of cycles, the class $\operatorname{SET}(\operatorname{CYC}(\mathcal Z))$ with series $\exp\log\frac1{1-z}=\frac1{1-z}$; so there are $n!$ of them, as the class $\operatorname{SEQ}(\mathcal Z)$ also says;
\item set partitions are sets of non-empty sets, the class $\operatorname{SET}(\operatorname{SET}_{\ge1}(\mathcal Z))$ with series $e^{e^z-1}$, which counts them by the Bell numbers;
\item ordered set partitions, or surjections onto $\{1,\dots,k\}$ for some $k$, are $\operatorname{SEQ}(\operatorname{SET}_{\ge1}(\mathcal Z))$: $\frac1{2-e^z}$ (Exercise~\ref{ex:fubini});
\item \emph{Cayley trees}, rooted trees on labelled vertices, are a root and a set of Cayley trees, $\mathcal T=\mathcal Z\star\operatorname{SET}(\mathcal T)$: $T(z)=z\,e^{T(z)}$;
\item maps $f:\{1,\dots,n\}\to\{1,\dots,n\}$ are sets of cycles of Cayley trees, $\operatorname{SET}(\operatorname{CYC}(\mathcal T))$: $\exp\log\frac1{1-T}=\frac1{1-T(z)}$.
\end{itemize}
The last decomposition deserves a proof. Draw the map $f$ as a graph with an arrow $x\to f(x)$ from each point. A point $x$ is \emph{cyclic} if $f^m(x)=x$ for some
$m\ge1$; the cyclic points are permuted by $f$, so they form cycles. Each point $x$ reaches a cyclic point after finitely many steps (the
sequence $x,f(x),f^2(x),\dots$ takes finitely many values, so repeats). Hang each non-cyclic $x$ below $f(x)$: the points whose first cyclic
point is $c$ form a tree rooted at $c$, the arrows pointing to the root. So $f$ is a set of cycles of rooted trees, and conversely such a
structure defines a map. The equation $T=ze^T$ has no closed solution, but the next section reads its coefficients exactly.

\section{Lagrange inversion}

The equations $G=\frac z{1-G}$ for plane trees and $T=ze^T$ for Cayley trees have the same shape: $T=z\varphi(T)$, ``a tree is a root and a
$\varphi$-structure of trees''. Lagrange found in 1770 how to invert such equations, for the needs of celestial mechanics. His formula is
a theorem about power series; here is a proof that it is really the cycle lemma, found by George Raney in 1960.

\begin{theorem}[Lagrange inversion]\label{thm:lagrange}
Let $\varphi(w)=\sum_i\varphi_iw^i\in\Q[[w]]$. There is a unique $T\in\Q[[z]]$ with $T(0)=0$ and $T=z\varphi(T)$, and for $n,k\ge1$,
\[
[z^n]\,T(z)^k=\frac kn\,[w^{n-k}]\,\varphi(w)^n .
\]
More generally, for every $H\in\Q[[w]]$ and $n\ge1$, $[z^n]H(T(z))=\frac1n[w^{n-1}]H'(w)\varphi(w)^n$.
\end{theorem}
\begin{proof}
\emph{Existence and uniqueness.} The coefficient of $z^n$ in $z\varphi(T)$ only involves the coefficients of $T$ below $n$, since $T^i$ has order at least $i$.
So $T=z\varphi(T)$ determines $t_1=\varphi_0$, $t_2,\dots$ one after the other. Each $t_n$, and each coefficient $[z^n]T^k$, is a polynomial in $\varphi_0,\dots,\varphi_{n-1}$ with
integer coefficients.

\emph{The case of natural numbers.} Suppose first that every $\varphi_i$ is a natural number. Let $\mathcal F$ be the class of plane trees in which each
vertex with $i$ children carries one of $\varphi_i$ colours, the size being the number of vertices. A tree is a root, a colour among $\varphi_i$, and a
sequence of $i$ trees, so $F=\sum_i\varphi_izF^i=z\varphi(F)$, and $F=T$ by uniqueness. Then $T^k$ is the OGF of the $k$-tuples of such trees. By
Proposition~\ref{prop:lukasiewicz} (whose proof treats codes of $k$ trees), listing the numbers of children of all vertices in preorder,
tree after tree, is a bijection from $k$-tuples of plane trees with $n$ vertices in all to the sequences $(c_1,\dots,c_n)\in\N^n$ such that
$h_j=k+\sum_{i\le j}(c_i-1)$ is $\ge1$ for $j<n$ and $h_n=0$. Colours ride along: the weight of a sequence is $\prod_i\varphi_{c_i}$, the number of ways to
colour the vertices.

As in the proof of Proposition~\ref{prop:lukasiewicz}, with $x_i=1-c_i$, the condition says that $(x_n,\dots,x_1)$ has positive partial sums, and its total
is $k$ exactly when $\sum_ic_i=n-k$. By the cycle lemma, every $(c_1,\dots,c_n)$ with $\sum c_i=n-k$ has exactly $k$ rotations that satisfy the condition.
Count the pairs (sequence, rotation satisfying the condition), each with the weight of the sequence, which does not change under
rotation: they have total weight $k\sum_c\prod_i\varphi_{c_i}=k\,[w^{n-k}]\varphi(w)^n$, and each valid sequence occurs in exactly $n$ pairs. Hence
$n\,[z^n]T^k=k\,[w^{n-k}]\varphi(w)^n$.

\emph{All coefficients.} For fixed $n$ and $k$, both sides are polynomials in $\varphi_0,\dots,\varphi_n$ with rational coefficients, and they agree whenever
these are natural numbers. A polynomial in $m$ variables that vanishes on $\N^m$ is zero (by induction on $m$: a polynomial in one variable
with infinitely many roots is zero, and the coefficients of the powers of the last variable vanish on $\N^{m-1}$). So the identity holds
for all $\varphi$.

\emph{The general form.} For $H=w^k$ with $k\ge1$ it is the identity just proved, since $\frac1n[w^{n-1}]kw^{k-1}\varphi^n=\frac kn[w^{n-k}]\varphi^n$; for $k=0$ both sides
vanish. By linearity it holds for polynomials, and for a series $H$ the coefficient of $z^n$ only involves $h_0,\dots,h_n$, since $T^k$ has order $\ge k$.
\end{proof}

\begin{example}[Cayley's formula, third proof]\label{ex:cayleylagrange}
For Cayley trees, $\varphi(w)=e^w$: the number of rooted trees on $n$ labelled vertices is
\[
\begin{aligned}{}&n!\,[z^n]T=\frac{n!}n\,[w^{n-1}]e^{nw}\eqbr=(n-1)!\,\frac{n^{n-1}}{(n-1)!}=n^{n-1}.\end{aligned}
\]
Each tree has $n$ choices of root, so there are $n^{n-2}$ trees, in agreement with Theorem~\ref{thm:prufer}.
\end{example}

For plane trees, $\varphi(w)=\frac1{1-w}$ and $[z^n]G=\frac1n[w^{n-1}](1-w)^{-n}=\frac1n\binom{2n-2}{n-1}=C_{n-1}$ (Tome~4, Proposition~2.1); for $r$-ary trees, $\varphi(w)=1+w^r$ gives the
Fuss--Catalan numbers of Exercise~\ref{ex:fusscatalan} in one line; for unary--binary trees, $\varphi(w)=1+w+w^2$ gives the Motzkin numbers as
coefficients of trinomials (Exercise~\ref{ex:motzkinlag}).

\begin{remark}
Lagrange's formula is the main tool for counting trees and maps, and it goes much further: Tutte counted planar maps with it in
the 1960s, and it is the starting point of the multivariate formulas of Good, used for branching processes with several types.
The classical analytic proof uses the residue theorem (Tome~4, Chapter~9); see Flajolet and Sedgewick, \emph{Analytic Combinatorics},
Appendix~A.6, or Stanley, \emph{Enumerative Combinatorics}, Volume~2, Section~5.4.
\end{remark}

\section{Parameters and averages}

The symbolic method counts more than objects. To study a parameter $\chi$ (the number of cycles of a permutation, the depth of the
vertices of a tree), use a second variable $u$ that records it: $F(z,u)=\sum_\alpha z^{|\alpha|}u^{\chi(\alpha)}$. Then $F(z,1)$ is the ordinary generating function, and
if all objects of size $n$ are equally likely, the mean of $\chi$ is
\[
\mathbb E_n[\chi]=\frac{[z^n]\,\partial_uF(z,1)}{[z^n]\,F(z,1)},
\]
since $\partial_uF(z,1)=\sum_\alpha\chi(\alpha)z^{|\alpha|}$; and the second derivative gives $\mathbb E_n[\chi(\chi-1)]$, hence the variance. The dictionary still applies when the
parameter is additive over the components: a marked component contributes $u$.

\begin{example}[cycles of a random permutation]\label{ex:cyclesbgf}
Marking each cycle with $u$, permutations are $\operatorname{SET}(u\operatorname{CYC}(\mathcal Z))$, so $F(z,u)=\exp\big(u\log\frac1{1-z}\big)=(1-z)^{-u}$. Its coefficients are
$[z^n](1-z)^{-u}=\frac{u(u+1)\cdots(u+n-1)}{n!}$, so the number of permutations of $n$ elements with $k$ cycles is the coefficient of $u^k$ in $u(u+1)\cdots(u+n-1)$, the
Stirling number $\stirlingi nk$ of Tome~4, Chapter~1. And $\partial_uF(z,1)=\frac1{1-z}\log\frac1{1-z}=\sum_nH_nz^n$, so a random permutation has $H_n\sim\ln n$ cycles on average
(Tome~0 found this mean by another route; Exercise~\ref{ex:cyclevar} gives the variance).
\end{example}

\begin{figure}[ht]
\centering
\fig{ch04-randomtree}
\caption{A uniformly random plane tree with $500$ vertices, drawn by the notebook (the root at the top, depth going down).
Its height is $40$, and the average depth of its vertices is $17.6$, of the order of $\frac12\sqrt{500\pi}\approx19.8$ predicted by Theorem~\ref{thm:pathlength}. A
random binary search tree with $500$ keys has average depth about $9.6$, growing like $2\ln n$ (Tome~4, Example~2.8 counts the same
quantity as the comparisons of quicksort): random trees in the uniform sense are tall and thin.}
\label{fig:randomtree}
\end{figure}

Now the depth of trees. The \emph{depth} of a vertex is its distance to the root, and the \emph{path length} of a tree is the sum of the depths of its
vertices. For a search tree, it is the total cost of finding every key; for a parse tree or a file system, the total length of all
paths.

\begin{theorem}[path length of plane trees]\label{thm:pathlength}
The sum of the path lengths of the $C_{n-1}$ plane trees with $n$ vertices is $\frac12\big(4^{n-1}-\binom{2n-2}{n-1}\big)$. Hence the depth of a uniformly random
vertex of a uniformly random plane tree with $n$ vertices has mean
\[
\begin{aligned}{}&\frac12\left(\frac{4^{n-1}}{\binom{2n-2}{n-1}}-1\right)\eqbr\sim\frac{\sqrt{\pi n}}2 .\end{aligned}
\]
\end{theorem}
\begin{proof}
Let $G(z)=\sum_tz^{|t|}$ and $P(z)=\sum_t\mathrm{pl}(t)z^{|t|}$, over plane trees $t$. A tree $t$ is a root with subtrees $s_1,\dots,s_k$, and every vertex of $s_i$ is one level
deeper in $t$ than in $s_i$, so $\mathrm{pl}(t)=\sum_i\big(\mathrm{pl}(s_i)+|s_i|\big)$. Summing over the trees with $k$ subtrees, and noting that $\sum_sz^{|s|}|s|=zG'(z)$, the
$i$-th subtree contributes $z\,G^{k-1}\big(P+zG'\big)$, so
\[
P=z\sum_{k\ge1}kG^{k-1}\big(P+zG'\big)=\frac z{(1-G)^2}\big(P+zG'\big).
\]
Now $G=\frac z{1-G}$ gives $\frac z{(1-G)^2}=\frac{G^2}z$ and $G^2=G-z$. So $zP=G^2P+zG^2G'$, that is, $P\,(2z-G)=zG^2G'$, where $2z-G=z-z^2-\cdots$ is $z$ times an
invertible series. Let $s=\sqrt{1-4z}$, the series with constant term $1$ whose square is $1-4z$. The root of $G^2-G+z=0$ with $G(0)=0$ is $G=\frac{1-s}2$, so
$G'=\frac1s$, $G^2=\frac{(1-s)^2}4$, and $2z-G=\frac{1-s^2}2-\frac{1-s}2=\frac{s(1-s)}2$. So
\[
\begin{aligned}{}&P(z)=\frac{z(1-s)}{2s^2}=\frac z2\Big(\frac1{1-4z}-\frac1{\sqrt{1-4z}}\Big),\end{aligned}
\]
and since $\frac1{\sqrt{1-4z}}=\sum_m\binom{2m}mz^m$ (Tome~4, Chapter~11), $[z^n]P=\frac12\big(4^{n-1}-\binom{2n-2}{n-1}\big)$. Dividing by the $nC_{n-1}=\binom{2n-2}{n-1}$ pairs (tree, vertex)
gives the mean, and $\binom{2m}m\sim\frac{4^m}{\sqrt{\pi m}}$ (Tome~4, Example~11.11) with $m=n-1$ gives its equivalent.
\end{proof}

For $n=3$: the path (depths $0,1,2$) and the cherry ($0,1,1$) have path lengths $3+2=5=\frac12(16-6)$. The answer, a square root, says that uniform
random trees are tall: unlike random binary search trees, which have logarithmic depth (Figure~\ref{fig:randomtree}). The height itself,
the largest depth, was studied by de Bruijn, Knuth and Rice in 1972: its mean is asymptotic to $\sqrt{\pi n}$, twice the mean depth.
David Aldous showed in 1991 that, rescaled by $\sqrt n$, the random tree converges to a random fractal, the \emph{Brownian continuum random tree},
whose contour is a Brownian excursion.

\section{Exercises}

\begin{exercise}[$\star$]\label{ex:compositions}
Find the OGF of the compositions of $n$ into parts $1$ and $2$, and of the compositions of $n$ into odd parts. Explain the surprise.
\end{exercise}

\begin{exercise}[$\star$]\label{ex:necklacecyc}
Derive the necklace formula $N(n,q)=\frac1n\sum_{d\mid n}\varphi(d)\,q^{n/d}$ of Tome~2, Example~3.8, from the cycle construction.
\end{exercise}

\begin{exercise}[$\star$]\label{ex:eulerodd}
Show with the dictionary that the partitions of $n$ into distinct parts are as many as the partitions of $n$ into odd parts (Euler, 1748).
\end{exercise}

\begin{exercise}[$\star$]\label{ex:mod3}
Build an automaton with three states reading binary words (most significant bit first) that accepts the representations of
multiples of $3$, and deduce that $\frac{2^n+2}3$ words of length $n$ are accepted for $n$ even, $\frac{2^n+1}3$ for $n$ odd. Check directly.
\end{exercise}

\begin{exercise}[$\star$]\label{ex:avoid}
Show that the binary words of length $n$ avoiding HH are counted by $F_{n+2}$, and those avoiding HT by $n+1$, using
Corollary~\ref{cor:penney}. Explain the second answer directly.
\end{exercise}

\begin{exercise}[$\star\star$]\label{ex:penney}
For each of the $8$ patterns $B=b_1b_2b_3$ of length $3$ over $\{\mathrm H,\mathrm T\}$, compute the odds with which $A=\bar b_2b_1b_2$ appears before $B$. Deduce that no pattern of
length $3$ is better than all the others, and verify the cycle of Figure~\ref{fig:penney}.
\end{exercise}

\begin{exercise}[$\star\star$]\label{ex:fubini}
The \emph{Fubini numbers} $f_n$ count the ordered set partitions of $\{1,\dots,n\}$ (rankings with ties of $n$ horses). Show that $\sum_nf_n\frac{z^n}{n!}=\frac1{2-e^z}$, compute
$f_1,\dots,f_5$, and use Tome~4, Theorem~11.13 to show that $f_n\sim\frac{n!}{2(\ln2)^{n+1}}$.
\end{exercise}

\begin{exercise}[$\star$]\label{ex:cyclevar}
Show that the number of cycles of a random permutation of $n$ elements has variance $H_n-H_n^{(2)}$, where $H_n^{(2)}=\sum_{k\le n}\frac1{k^2}$.
\end{exercise}

\begin{exercise}[$\star\star$]\label{ex:maps}
Let $T=ze^T$ be the EGF of Cayley trees. Show that $zT'=\frac T{1-T}$, so that the EGF of maps is $\frac1{1-T}=1+zT'$. Since there are $n^n$ maps from $\{1,\dots,n\}$ to itself,
deduce a fourth proof that there are $n^{n-1}$ Cayley trees, without Lagrange inversion.
\end{exercise}

\begin{exercise}[$\star\star$]\label{ex:forests}
Show that the number of rooted forests on $\{1,\dots,n\}$ with $k$ trees is $\binom{n-1}{k-1}n^{n-k}$.
\end{exercise}

\begin{exercise}[$\star$]\label{ex:motzkinlag}
Unary--binary plane trees have vertices with $0$, $1$ or $2$ children. Show that there are $\frac1n[w^{n-1}](1+w+w^2)^n$ of them with $n$ vertices, and compute
the first six values.
\end{exercise}

\begin{exercise}[$\star\star$]\label{ex:leaves}
Mark the leaves of plane trees with $u$. Show that the number of plane trees with $n$ vertices and $j$ leaves is the \emph{Narayana number}
$\frac1n\binom nj\binom{n-2}{j-1}$ for $n\ge2$, and that a random plane tree with $n\ge2$ vertices has exactly $\frac n2$ leaves on average.
\end{exercise}

\section{Hints}

\begin{hint}{ex:compositions}
$\operatorname{SEQ}(\mathcal Z+\mathcal Z^2)$, and $\operatorname{SEQ}$ of $\mathcal Z+\mathcal Z^3+\mathcal Z^5+\cdots$.
\end{hint}
\begin{hint}{ex:necklacecyc}
Extract $[z^n]$ from $\sum_d\frac{\varphi(d)}d\log\frac1{1-qz^d}$.
\end{hint}
\begin{hint}{ex:eulerodd}
$1+z^n=\frac{1-z^{2n}}{1-z^n}$.
\end{hint}
\begin{hint}{ex:mod3}
The state is the remainder of the number read so far; reading a bit $b$ sends $r$ to $2r+b$ mod $3$.
\end{hint}
\begin{hint}{ex:avoid}
Find the correlation polynomials: $c_{\mathrm{HH,HH}}=1+z$ and $c_{\mathrm{HT,HT}}=1$.
\end{hint}
\begin{hint}{ex:penney}
Compute the four Conway numbers for each pair; the overlaps of length $j=1,2,3$ are worth $2,4,8$.
\end{hint}
\begin{hint}{ex:fubini}
$\operatorname{SEQ}(\operatorname{SET}_{\ge1}(\mathcal Z))$. For the asymptotics, the smallest zero of $2-e^z$ is $\ln2$; substitute $z=w\ln2$.
\end{hint}
\begin{hint}{ex:cyclevar}
$\partial_u^2(1-z)^{-u}$ at $u=1$ is $\frac1{1-z}\log^2\frac1{1-z}$.
\end{hint}
\begin{hint}{ex:maps}
Differentiate $T=ze^T$.
\end{hint}
\begin{hint}{ex:forests}
A forest with $k$ trees is $\operatorname{SET}_k(\mathcal T)$, with EGF $\frac{T^k}{k!}$.
\end{hint}
\begin{hint}{ex:motzkinlag}
$\varphi(w)=1+w+w^2$.
\end{hint}
\begin{hint}{ex:leaves}
$G=z\big(u-1+\frac1{1-G}\big)$: Lagrange with $\varphi(w)=u-1+\frac1{1-w}$, and the binomial theorem in $u$.
\end{hint}

\section{Solutions}

\begin{solution}{ex:compositions}
\idea Translate, then compare the two rational functions.
\writeup Parts $1$ and $2$: $\frac1{1-z-z^2}$, so $F_{n+1}$ compositions of $n$. Odd parts: the parts form $\frac z{1-z^2}$, and $\frac1{1-z/(1-z^2)}=\frac{1-z^2}{1-z-z^2}$, whose coefficient of
$z^n$ is $F_{n+1}-F_{n-1}=F_n$ for $n\ge1$. So the compositions of $n$ into odd parts are as many as the compositions of $n-1$ into parts $1$ and $2$:
for $n=5$, the five compositions $5$, $3+1+1$, $1+3+1$, $1+1+3$, $1+1+1+1+1$ match the five compositions of $4$ into $1$s and $2$s. Finding a
bijection that explains it is a pleasant puzzle.
\end{solution}

\begin{solution}{ex:necklacecyc}
\idea The coefficient of $z^n$ in $\log\frac1{1-qz^d}$ is easy: it is non-zero only when $d$ divides $n$.
\writeup $\log\frac1{1-qz^d}=\sum_m\frac{q^mz^{dm}}m$, so $[z^n]$ of it is $\frac{q^{n/d}}{n/d}$ if $d\mid n$ and $0$ otherwise. Hence the number of necklaces of length $n$, which are
the cycles of $\operatorname{CYC}(q\mathcal Z)$ of size $n$, is $\sum_{d\mid n}\frac{\varphi(d)}d\cdot\frac d n\,q^{n/d}=\frac1n\sum_{d\mid n}\varphi(d)q^{n/d}$.
\end{solution}

\begin{solution}{ex:eulerodd}
\idea Write both generating functions as products and cancel.
\writeup Distinct parts: $\operatorname{PSET}(\operatorname{SEQ}_{\ge1}(\mathcal Z))\mapsto\prod_n(1+z^n)=\prod_n\frac{1-z^{2n}}{1-z^n}$. In the last product, every factor $1-z^{2n}$ of the numerator cancels the
factor $1-z^m$, $m=2n$, of the denominator, leaving $\prod_{m\text{ odd}}\frac1{1-z^m}$, the OGF of $\operatorname{MSET}$ of the odd integers. (The cancellation is legitimate in
$\Q[[z]]$: truncate both products to the factors with exponent at most $N$ and compare the coefficients up to $z^N$.)
\end{solution}

\begin{solution}{ex:mod3}
\idea Reading a binary number from the left, the value doubles and adds the new bit, so its remainder modulo $3$ is all one needs to
remember.
\writeup States $0,1,2$ (the remainder of the number read so far), start and accept at $0$, $\delta(r,b)=2r+b\bmod3$: so $0\to\{0,1\}$, $1\to\{2,0\}$,
$2\to\{1,2\}$. With $L_r$ the OGF of the words leading from $r$ to $0$, Theorem~\ref{thm:automaton} gives
\[
\begin{gathered}L_0=1+z(L_0+L_1),\qq L_1=z(L_2+L_0),\qq L_2=z(L_1+L_2).\end{gathered}
\]
The last equation gives $L_2=\frac{zL_1}{1-z}$, then the second $L_1=\frac{z(1-z)}{1-z-z^2}L_0$, and the first
\[
\begin{aligned}{}&L_0=\frac{1-z-z^2}{(1-z)(1+z)(1-2z)}\eqbr=\frac13\Big(\frac1{1-2z}+\frac1{1-z}+\frac1{1-z^2}\Big),\end{aligned}
\]
by partial fractions (check at $z=0$ and by clearing denominators). The coefficient of $z^n$ is $\frac{2^n+2}3$ for $n$ even and $\frac{2^n+1}3$ for $n$ odd. Directly: the
words of length $n$ are the integers $0\le x<2^n$ written with leading zeros, and the multiples of $3$ among them number $\lceil2^n/3\rceil$, which is the
same, as $2^n\equiv1$ or $2\pmod3$ according to the parity of $n$.
\end{solution}

\begin{solution}{ex:avoid}
\idea Read off the overlaps and substitute in Corollary~\ref{cor:penney}(1) with $q=2$, $k=2$.
\writeup HH overlaps itself in $j=1$ and $j=2$, so $c=z+1$, and $S=\frac{1+z}{z^2+(1-2z)(1+z)}=\frac{1+z}{1-z-z^2}$, whose coefficients are $F_{n+1}+F_n=F_{n+2}$. HT overlaps itself
only in $j=2$, so $c=1$ and $S=\frac1{z^2+1-2z}=\frac1{(1-z)^2}$, with coefficients $n+1$. Directly: a word avoiding HT never has an H followed by a T, so it is
$\mathrm T^i\mathrm H^{n-i}$, for $i=0,\dots,n$.
\end{solution}

\begin{solution}{ex:penney}
\idea For each pair, list the overlaps (suffix of one equal to prefix of the other) and apply Conway's formula; the overlaps
of $A=\bar b_2b_1b_2$ with $B$ are favourable because $A$ ends with the beginning $b_1b_2$ of $B$.
\writeup The odds of $A$ are $(BB-BA):(AA-AB)$. For $B=\mathrm{HHH}$, $A=\mathrm{THH}$: $BB=2+4+8=14$, $BA=0$, $AA=8$, $AB=2+4=6$: odds $14:2=7:1$. For $B=\mathrm{HHT}$,
$A=\mathrm{THH}$: $BB=8$, $BA=2$, $AA=8$, $AB=6$: $6:2=3:1$. For $B=\mathrm{HTH}$, $A=\mathrm{HHT}$: $BB=8+2=10$, $BA=2$ (H), $AA=8$, $AB=4$ (HT): $8:4=2:1$. For $B=\mathrm{HTT}$,
$A=\mathrm{HHT}$: $BB=8$, $BA=0$, $AA=8$, $AB=4$ (HT): $8:4=2:1$. The other four follow by exchanging H and T. So every pattern is beaten by another one,
with odds at least $2:1$, and none can be the best. The cycle: THH beats HHT ($3:1$), TTH beats THH ($2:1$, by the symmetric of the case $B=\mathrm{HTT}$),
HTT beats TTH ($3:1$, symmetric of $B=\mathrm{HHT}$), HHT beats HTT ($2:1$).
\end{solution}

\begin{solution}{ex:fubini}
\idea A ranking with ties is a sequence of non-empty blocks of horses; for the asymptotics, the dominant singularity is a simple
pole.
\writeup $\operatorname{SEQ}(\operatorname{SET}_{\ge1}(\mathcal Z))\mapsto\frac1{1-(e^z-1)}=\frac1{2-e^z}$. Expanding, $f_1=1$, $f_2=3$, $f_3=13$, $f_4=75$, $f_5=541$ (or by the recurrence $f_n=\sum_{k\ge1}\binom nkf_{n-k}$, which
says that the first block has $k$ horses). The zeros of $2-e^z$ are $\ln2+2\pi ik$, $k\in\Z$; the one closest to $0$ is $\ln2$, and the others have modulus
$>2\pi>\ln 2$. Put $G(w)=\frac1{2-e^{w\ln2}}=\frac{h(w)}{1-w}$ with $h(w)=\frac{1-w}{2-2^w}$. Then $h$ is analytic for $|w|<\frac{|\ln2+2\pi i|}{\ln2}$, a radius $>1$, with $h(1)=\frac1{2\ln2}$
(by l'H\^opital, the derivative of $2-2^w$ at $1$ being $-2\ln2$). Theorem~11.13 of Tome~4 with $\alpha=1$ gives $[w^n]G\to\frac1{2\ln2}$, and $[z^n]\frac1{2-e^z}=(\ln2)^{-n}[w^n]G$, so
$f_n=n!\,[z^n]\frac1{2-e^z}\sim\frac{n!}{2(\ln2)^{n+1}}$. Already for $n=5$, $\frac{120}{2(\ln2)^6}\approx541.0$.
\end{solution}

\begin{solution}{ex:cyclevar}
\idea Two derivatives of $(1-z)^{-u}=e^{uL}$, with $L=\log\frac1{1-z}$, give the factorial moments.
\writeup $\partial_u^2e^{uL}=L^2e^{uL}$, so $\mathbb E_n[C(C-1)]=[z^n]\frac{L^2}{1-z}$, the sum of the coefficients of $z^m$, $m\le n$, in $L^2$. That coefficient is
$\sum_{i+j=m}\frac1{ij}=\frac1m\sum_{i+j=m}\big(\frac1i+\frac1j\big)=\frac2mH_{m-1}$. And $\sum_{m\le n}\frac2mH_{m-1}=2\sum_{i<m\le n}\frac1{im}=H_n^2-H_n^{(2)}$, separating in $H_n^2=\sum_{i,m\le n}\frac1{im}$ the
terms $i=m$ from the others. Hence $\operatorname{Var}C=\mathbb E[C(C-1)]+\mathbb E[C]-\mathbb E[C]^2=H_n^2-H_n^{(2)}+H_n-H_n^2=H_n-H_n^{(2)}$.
The variance is about $\ln n$, like the mean: the number of cycles is concentrated, and Goncharov (1944) proved that it is asymptotically
normal.
\end{solution}

\begin{solution}{ex:maps}
\idea Differentiating the functional equation expresses $T'$ through $T$; then the two decompositions of maps give the same series.
\writeup By the chain rule, $T'=e^T+ze^TT'=\frac Tz+TT'$, so $T'(1-T)=\frac Tz$ and $zT'=\frac T{1-T}$. Hence $\frac1{1-T}=1+\frac T{1-T}=1+zT'$. The coefficient of $\frac{z^n}{n!}$ in
$\frac1{1-T}$ is the number of maps, $n^n$; in $1+zT'$ it is $n\,t_n$ for $n\ge1$, where $t_n$ is the number of Cayley trees (since $zT'=\sum_nnt_n\frac{z^n}{n!}$). So $nt_n=n^n$ and
$t_n=n^{n-1}$. This is, in disguise, Joyal's bijective proof (1981): a map is a rooted tree with a second marked vertex.
\end{solution}

\begin{solution}{ex:forests}
\idea A forest is a set of $k$ Cayley trees; Lagrange gives the coefficients of $T^k$.
\writeup The EGF is $\frac{T^k}{k!}$, and Theorem~\ref{thm:lagrange} with $\varphi=e^w$ gives $[z^n]T^k=\frac kn[w^{n-k}]e^{nw}=\frac kn\cdot\frac{n^{n-k}}{(n-k)!}$. So the number of forests is
$\frac{n!}{k!}\cdot\frac kn\cdot\frac{n^{n-k}}{(n-k)!}=\frac{(n-1)!}{(k-1)!\,(n-k)!}n^{n-k}=\binom{n-1}{k-1}n^{n-k}$. Checks: $k=1$ gives $n^{n-1}$, $k=n$ gives $1$ (all singletons).
\end{solution}

\begin{solution}{ex:motzkinlag}
\idea A unary--binary tree is a root with $0$, $1$ or $2$ ordered subtrees: $M=z(1+M+M^2)$.
\writeup Theorem~\ref{thm:lagrange} with $\varphi(w)=1+w+w^2$ gives $[z^n]M=\frac1n[w^{n-1}](1+w+w^2)^n$. For $n=1,\dots,6$: $1$, $\frac12\cdot2=1$, $\frac13\cdot6=2$, $\frac14\cdot16=4$, $\frac15\cdot45=9$,
$\frac16\cdot126=21$: the Motzkin numbers $1,1,2,4,9,21$ (Tome~4, Exercise~2.6 met them through another equation).
\end{solution}

\begin{solution}{ex:leaves}
\idea A leaf is a vertex with $0$ children: in the Lagrange formula, giving weight $u$ to the term $\varphi_0$ marks the leaves.
\writeup A tree is a root with $i$ subtrees, and the root is a leaf if $i=0$: $G=z\big(u+G+G^2+\cdots\big)=z\varphi(G)$ with $\varphi(w)=u-1+\frac1{1-w}=u+\frac w{1-w}$. By
Theorem~\ref{thm:lagrange}, applied for each rational value of $u$ (both sides below are polynomials in $u$, equal for all rational $u$, hence equal),
\[
\begin{aligned}{}&[z^n]G=\frac1n[w^{n-1}]\Big(u+\frac w{1-w}\Big)^n\eqbr=\frac1n\sum_j\binom nju^j[w^{n-1}]\frac{w^{n-j}}{(1-w)^{n-j}}\eqbr=\frac1n\sum_j\binom nju^j\binom{n-2}{j-1},\end{aligned}
\]
using $[w^{j-1}](1-w)^{-(n-j)}=\binom{n-2}{j-1}$ (Tome~4, Proposition~2.1). For the mean, $\sum_jj\frac1n\binom nj\binom{n-2}{j-1}=\sum_j\binom{n-1}{j-1}\binom{n-2}{j-1}=\binom{2n-3}{n-1}$ (Vandermonde), and
dividing by $C_{n-1}=\frac1n\binom{2n-2}{n-1}$ gives $\frac{n\binom{2n-3}{n-1}}{\binom{2n-2}{n-1}}=n\cdot\frac{n-1}{2n-2}=\frac n2$.
\end{solution}

\section*{Sources and further reading}

\begin{itemize}
\item W.~Penney, \emph{Problem 95: Penney-Ante}, Journal of Recreational Mathematics 2 (1969), 241; M.~Gardner, \emph{Mathematical
Games}, Scientific American, October 1974.
\item L.~J.~Guibas and A.~M.~Odlyzko, \emph{String overlaps, pattern matching, and nontransitive games}, Journal of Combinatorial
Theory A 30 (1981), 183--208. A research paper at the level of this chapter: the equations of Theorem~\ref{thm:guibas} for any
number of patterns, the asymptotics of the waiting times, and the game.
\item P.~Flajolet and R.~Sedgewick, \emph{Analytic Combinatorics}, Cambridge University Press, 2009 (free online). The symbolic
method, labelled and unlabelled, and its analytic half.
\item R.~Sedgewick and P.~Flajolet, \emph{An Introduction to the Analysis of Algorithms}, 2nd ed., Addison-Wesley, 2013.
\item D.~E.~Knuth, \emph{The Art of Computer Programming}, Volume~1, 3rd ed., Addison-Wesley, 1997 (trees, path length, Lagrange
inversion in Section~2.3.4.4).
\item G.~N.~Raney, \emph{Functional composition patterns and power series reversion}, Transactions of the AMS 94 (1960),
441--451. The combinatorial proof of Lagrange inversion.
\item N.~Chomsky and M.-P.~Sch\"utzenberger, \emph{The algebraic theory of context-free languages}, in \emph{Computer Programming
and Formal Systems}, North-Holland, 1963, 118--161.
\item N.~G.~de Bruijn, D.~E.~Knuth and S.~O.~Rice, \emph{The average height of planted plane trees}, in \emph{Graph Theory and
Computing}, Academic Press, 1972, 15--22; D.~Aldous, \emph{The continuum random tree I}, Annals of Probability 19 (1991), 1--28.
\item A.~Joyal, \emph{Une th\'eorie combinatoire des s\'eries formelles}, Advances in Mathematics 42 (1981), 1--82.
\item H.~S.~Wilf, \emph{generatingfunctionology}, 3rd ed., A K Peters, 2006 (free online).
\end{itemize}
